% NOTE_KOIDE_DERIVATION_2026-05-24.tex
% Partial Route-A geometric result for the Koide identity K = 2/3
% Author: Kévin Rémondière, 2026-05-24
% License: CC-BY 4.0
\documentclass[11pt,a4paper]{article}

\usepackage[utf8]{inputenc}
\usepackage[T1]{fontenc}
\usepackage{amsmath, amssymb, amsthm}
\usepackage{geometry}
\geometry{margin=2.4cm}
\usepackage{hyperref}
\hypersetup{colorlinks=true, linkcolor=blue, citecolor=blue, urlcolor=blue}
\usepackage{authblk}
\usepackage{graphicx}
\usepackage{booktabs}

\newtheorem{theorem}{Theorem}
\newtheorem{lemma}[theorem]{Lemma}
\newtheorem{proposition}[theorem]{Proposition}
\newtheorem{corollary}[theorem]{Corollary}
\theoremstyle{remark}
\newtheorem*{remark}{Remark}

\title{The Koide identity $K = 2/3$ as a variance--rank constraint on the lepton sqrt--mass simplex : a partial geometric result}
\author{K\'evin R\'emond\`iere\thanks{Independent researcher, Oloron-Sainte-Marie, 64400 France. ORCID 0009-0008-2443-7166. Email \texttt{kevin.remondiere@gmail.com}}}
\date{2026-05-24}

\begin{document}

\maketitle

\begin{abstract}
We present a partial geometric result for the Koide identity $K \equiv (m_e + m_\mu + m_\tau)/(\sqrt{m_e}+\sqrt{m_\mu}+\sqrt{m_\tau})^2 = 2/3$. Setting $u_i := \sqrt{m_i}/\sum_j\sqrt{m_j}$ (normalised square--root masses), the Koide identity is equivalent to the variance constraint $\mathrm{Var}(u_i) = 1/N^2 = 1/9$ for $N = 3$, equivalently to the angle constraint $\theta(\vec u, (1,1,1)/\sqrt 3) = \pi/4 = 45^\circ$. We document with full precision that the PDG 2024 lepton masses satisfy this constraint at $0.91\sigma$, that the constraint is a sharp non-trivial condition on the Yukawa hierarchy ($\sim 5 \times 10^{-3}$ probability density under uniform sampling on the 2--simplex), and that the value $N = 3$ coincides with $|\Phi^+(\mathrm{SU}(3))|$, the cardinality of positive roots of the colour gauge group $A_2$. We are explicit that this is a \emph{geometric rephrasing}, not a first--principles \emph{derivation}, and we state the two open lemmas that would close the gap to a full derivation.
\end{abstract}

\section{The Koide identity and its variance form}

For three positive real numbers $m_1, m_2, m_3$ define
\begin{equation}
K(m_1, m_2, m_3) := \frac{m_1 + m_2 + m_3}{\big(\sqrt{m_1} + \sqrt{m_2} + \sqrt{m_3}\big)^2}.
\label{eq:Koide-def}
\end{equation}
Elementary inequalities give $K \in [1/3, 1]$ with $K = 1/3$ at $m_1 = m_2 = m_3$ and $K \to 1$ as one mass dominates. The Koide observation \cite{Koide1983} is that the empirical value at the charged--lepton pole masses lies very close to $2/3$ :
\begin{equation}
K(m_e, m_\mu, m_\tau)\big|_{\mathrm{PDG 2024}} = 0.66666051 \pm 0.00000677,
\label{eq:K-PDG}
\end{equation}
i.e. $0.91\sigma$ from the rational $2/3 = 0.6666\overline{6}$, with $\sigma_K = 6.77 \times 10^{-6}$ propagated from $\sigma_{m_\tau} = 0.12~\mathrm{MeV}$ via $|\partial K/\partial m_\tau| \approx 5.6 \times 10^{-5}~\mathrm{MeV}^{-1}$.

Let $u_i := \sqrt{m_i}$. Then $K(\vec m) = \|\vec u\|^2 / \langle \vec u, \vec{\mathbf{1}}\rangle^2$ where $\vec{\mathbf{1}} = (1, 1, 1)$, so $K$ is invariant under positive scaling. Normalise by $S := \sum_i u_i$ and set $\tilde u_i := u_i / S$ ; then
\begin{equation}
K(\vec m) = \sum_{i=1}^3 \tilde u_i^2, \qquad \sum_i \tilde u_i = 1, \qquad \tilde u_i \ge 0.
\label{eq:K-as-sum-squares}
\end{equation}
The variance of the normalised components $\tilde u_i$ around their common mean $\bar{\tilde u} = 1/3$ is
\begin{equation}
\mathrm{Var}(\tilde u_i) = \frac{1}{3}\sum_i \tilde u_i^2 - \frac{1}{9} = \frac{K}{3} - \frac{1}{9},
\end{equation}
so $K = 2/3 \iff \mathrm{Var}(\tilde u_i) = 1/9$.

\begin{proposition}[Variance--rank reformulation]
\label{prop:variance-rank}
For three positive masses $m_1, m_2, m_3$, the Koide identity $K = 2/N$ with $N = 3$ is equivalent to
\begin{equation}
\boxed{\;\mathrm{Var}\big(\tilde u_1, \tilde u_2, \tilde u_3\big) = \frac{1}{N^2} = \frac{1}{9}, \quad N = 3.\;}
\label{eq:variance-rank}
\end{equation}
\end{proposition}

\begin{remark}
The variance is computed at the discrete probability $1/3$ on each of the three index values, not in any continuous--probability sense. The right--hand side $1/N^2$ uses the dimension $N$ of the relevant simplex ($N = 3$ for three generations) ; equivalently, the standard deviation $\sigma(\tilde u_i) = 1/N = 1/3$.
\end{remark}

\section{Foot's $45^\circ$ angle in this language}

Set $\vec n := \vec{\mathbf{1}}/\sqrt 3$, the unit vector along the diagonal of $\mathbb{R}^3$. For any $\vec u \in \mathbb{R}^3_+$ define $\cos\theta := \langle \vec u, \vec n\rangle / \|\vec u\|$. Direct algebra gives
\begin{equation}
K = \frac{\|\vec u\|^2}{\langle \vec u, \vec{\mathbf{1}}\rangle^2} = \frac{\|\vec u\|^2}{3 \, \|\vec u\|^2 \cos^2\theta} = \frac{1}{3\cos^2\theta},
\end{equation}
so $K = 2/3 \iff \cos^2\theta = 1/2 \iff \theta = \pi/4 = 45^\circ$. This is Foot's geometric interpretation \cite{Foot1994}.

\begin{proposition}[Empirical confirmation at PDG 2024]
\label{prop:empirical}
With PDG 2024 charged--lepton pole masses, the angle $\theta_{\mathrm{emp}} = \arccos\big(\langle \vec u_{\mathrm{emp}}, \vec n\rangle / \|\vec u_{\mathrm{emp}}\|\big)$ between the empirical lepton sqrt--mass vector and the diagonal is
\begin{equation}
\theta_{\mathrm{emp}} = 44.99974^\circ \pm 0.00040^\circ,
\end{equation}
consistent with $45^\circ$ at $0.91\sigma$ in $K$--space, equivalently $\sim 6 \times 10^{-6}$ relative deviation in $\cos^2\theta$.
\end{proposition}

\section{Statistical significance under uniform priors}

\begin{proposition}[Hierarchy as a non--trivial constraint]
\label{prop:significance}
Let $\vec u$ be uniformly distributed on the positive 2--sphere $\{\vec u \in \mathbb{R}^3_+ : \|\vec u\| = 1\}$ (equivalently $\vec m \sim |\vec u|^2$ uniform in direction). Then the induced distribution of $K(\vec u) = 1/(3\cos^2\theta)$ satisfies
\begin{align}
\mathbb{E}[K] &\approx 0.461, & \mathrm{median}\,K &\approx 0.436, \\
\Pr\big(K \in [0.66, 0.67]\big) &\approx 5.3 \times 10^{-3}, & \Pr\big(|K - 2/3| < 10^{-5}\big) &\approx 2 \times 10^{-5}.
\end{align}
\end{proposition}

\begin{proof}[Sketch]
Numerical estimate over $10^5$ uniform samples on the positive 2--sphere ; see the companion calculation script \texttt{koide\_test\_all\_families.py}. The probability of $K$ landing within $10^{-5}$ of $2/3$ by chance is approximately $2 \times 10^{-5}$, four orders of magnitude smaller than the empirical four--significant--figure agreement. The lepton hierarchy is therefore a sharp non--trivial constraint.
\end{proof}

\section{Coincidence $N = |\Phi^+(\mathrm{SU}(3))|$ and the structural overlay}

For the colour gauge group $\mathrm{SU}(N)$ with root system $A_{N-1}$, the number of positive roots is $|\Phi^+(\mathrm{SU}(N))| = N(N-1)/2$. The equality
\begin{equation}
N = |\Phi^+(\mathrm{SU}(N))| \iff N - 1 = 2 \iff N = 3
\label{eq:N-equals-Phi}
\end{equation}
holds \emph{uniquely} at $N = 3$. At $N = 3$, the rephrasing of Proposition~\ref{prop:variance-rank} therefore admits the dual reading
\begin{equation}
K = \frac{2}{N} = \frac{2}{|\Phi^+(\mathrm{SU}(3))|} = \frac{2}{3}.
\end{equation}
A companion log--Sobolev framework for compact--Lie--group Wilson lattice gauge theory \cite{Remondiere2026b} introduces the Lie--algebraic saturation coefficient
\begin{equation}
\kappa(G) := \frac{1}{2|\Phi^+(G)|},
\end{equation}
giving $\kappa(\mathrm{SU}(3)) = 1/6$ (kernel--verified in Lean 4). The empirical--structural identification is then the one--line identity
\begin{equation}
K_{\mathrm{Koide}} = 4\kappa(\mathrm{SU}(3)) = \frac{2}{|\Phi^+(\mathrm{SU}(3))|} = \frac{2}{3}
\label{eq:K-4kappa}
\end{equation}
at the four--significant--figure level of present experimental precision.

\section{Honest scope : this is a rephrasing, not a derivation}

We are explicit about what this note does and does not establish.

\begin{itemize}
\item \textbf{Established :} the geometric rephrasing $K = 2/3 \iff \sigma(\tilde u_i) = 1/N \iff \theta = \pi/4$ (Propositions~\ref{prop:variance-rank}, \ref{prop:empirical}) ; the statistical significance of the empirical agreement (Proposition~\ref{prop:significance}) ; the coincidence $N = |\Phi^+(\mathrm{SU}(3))|$ uniquely at $N = 3$ (Eq.~\ref{eq:N-equals-Phi}).

\item \textbf{Not established :} a \emph{causal} link between the variance constraint Eq.~\ref{eq:variance-rank} and the Lie--algebraic coefficient $\kappa(\mathrm{SU}(3)) = 1/6$. The two are numerically compatible via Eq.~\ref{eq:K-4kappa}, but no Lagrangian--level mechanism is known by which the colour--sector $\kappa$ would force the lepton--sector variance to take the value $1/9$.
\end{itemize}

\section{Two open lemmas}

The companion long--form report \cite{Remondiere2026deriv} identifies two specific open lemmas whose proof would close the gap to a full derivation :

\paragraph{Lemma 1 (Variance--rank from flavour symmetry breaking).} If $\mathrm{SU}(N)_{\mathrm{flavour}}$ is spontaneously broken to the trivial subgroup with vacuum alignment in the equal--power direction $\langle \phi \rangle = (\sqrt{m_1}, \ldots, \sqrt{m_N})$, then the variance constraint $\mathrm{Var}(\tilde u_i) = 1/N^2$ is equivalent to the requirement that the singlet $\mathrm{U}(1)$ component of $\langle \phi \rangle$ carries exactly half the norm--squared. \emph{What needs proving} : that this equal--power vacuum alignment is the unique energy minimum of a natural flavour--Higgs potential.

\paragraph{Lemma 2 (Higgs as radial mode of $\mathcal{A}/\mathcal{G}$).} Let $\mathcal{A}/\mathcal{G}$ be the zeta--regularised gauge orbit space of pure $\mathrm{SU}(3)$ Yang--Mills in four spacetime dimensions \cite{Singer1981, YangHartnoll2018}. \emph{What needs proving} : that the Higgs field is the radial mode of $\mathcal{A}/\mathcal{G}$, that its effective potential admits a minimum $r_* > 0$, and that the coupling of lepton singlets to $\mathcal{A}/\mathcal{G}$ via this radial mode produces a Yukawa matrix with $K = 2/3$ in equilibrium.

If both lemmas were proved, the chain $\kappa(\mathrm{SU}(3)) = 1/6 \Rightarrow$ Higgs radial mode $\Rightarrow$ flavour vacuum alignment $\Rightarrow$ $K = 2/3$ would constitute a first--principles derivation. Neither lemma is currently available ; both are research programmes.

\section*{Acknowledgments and COPE--compliant LLM disclosure}

The author thanks the Particle Data Group for the precision charged--lepton pole masses, and prior Koide--formula investigators (Y. Koide, R. Foot, A. Rivero, Y. Sumino, J. Kocik) for the conceptual scaffolding. In accordance with COPE recommendations on AI--assisted research \cite{COPE2023}, large--language--model assistance (Anthropic Claude, Opus 4 family, 2026--05) was used for adversarial reference cross--checking against the live arXiv API and for numerical verification of the variance and angle identities. All intellectual content, including the variance--rank reformulation of Proposition~\ref{prop:variance-rank}, the explicit honest scope statement of \S 5, and the two open lemmas of \S 6, is owned by the human author. The LLM is not an author.

\begin{thebibliography}{99}

\bibitem{Koide1983} Y. Koide, \emph{New view of quark and lepton mass hierarchy}, Phys. Rev. D \textbf{28}, 252 (1983).

\bibitem{Foot1994} R. Foot, \emph{A note on Koide's lepton mass relation}, arXiv:hep--ph/9402242 (1994).

\bibitem{Singer1981} I.M. Singer, \emph{The geometry of the orbit space for non--Abelian gauge theories}, Phys. Scripta \textbf{24}, 817 (1981).

\bibitem{YangHartnoll2018} H. Yang and S. Hartnoll, \emph{Orbit Space Curvature as a Source of Mass in Quantum Gauge Theory}, arXiv:1809.06318 (2018).

\bibitem{Sumino2009a} Y. Sumino, \emph{Family Gauge Symmetry and Koide's Mass Formula}, Phys. Lett. B \textbf{671}, 477 (2009), arXiv:0812.2090.

\bibitem{Sumino2009b} Y. Sumino, \emph{Family Gauge Symmetry as an Origin of Koide's Mass Formula and Charged Lepton Spectrum}, JHEP \textbf{0905}, 075 (2009), arXiv:0812.2103.

\bibitem{Remondiere2026b} K. R\'emond\`iere, \emph{Synthesis v2 : log--Sobolev framework for compact--Lie--group gauge theory} (with Lean 4 certification of $\kappa(\mathrm{SU}(3)) = 1/6$), manuscript and data archive, 2026 ; Zenodo concept DOI 10.5281/zenodo.19686398.

\bibitem{Remondiere2026deriv} K. R\'emond\`iere, \emph{OP--KOIDE--DERIVATION : Attempt to derive $K = 4\kappa = 2/3$ from first principles}, companion long--form report, 2026--05--24.

\bibitem{PDG2024} R.L. Workman et al. (Particle Data Group), \emph{Review of Particle Physics}, Prog. Theor. Exp. Phys. \textbf{2024}, 083C01 (2024).

\bibitem{COPE2023} Committee on Publication Ethics, \emph{Authorship and AI tools}, COPE position statement (2023).

\end{thebibliography}

\end{document}
